\section{Related Work}
\label{sec:relatedwork}
% \tom{

We first review prior works for next best view computation. We then discuss depth estimation literature, from which we borrow techniques to avoid the need for depth acquisition.

\subsection{Next Best View (NBV)}
Approaches to NBV can be broadly split into two groups based on the scene representation. On the one hand, volumetric methods represent the scene as voxels used as inputs of traditional optimization schemes~\cite{potthast-14-probabilistic-framework, vasquez-14-volumetric-nbv, charrow-15-information, bissmarck-15-efficient-nbv, vasquez-17-view-state-planning, daudelin-17-adaptable-probabilistic, cieslewski-17-ICIRS, wang-20-efficient-autonomous} or more recently, neural networks~\cite{hepp-corr18-learntoscore, mendoza-prl2020, vasquez-21-3DCNN}. 
On the other hand, surface-based approaches~\cite{chen-smc05-visionsensorplanning, kriegel-11-surface-based, kriegel-12-next-best-scan, lee-tra20-mechanical-parts, zeng-icirs20-pcnbv} operate on dense point clouds representing the surface as computed by the depth sensor. Although modeling surfaces allows to preserve highly-detailed geometries, it does not scale well to complex scenes involving large point clouds, thus limiting their applicability to synthetic settings of isolated centered objects with cameras lying on a sphere. The closest work to ours is Guédon~\etal~\cite{guedon-nips22-scone-surface-coverage} which proposes an hybrid approach called SCONE that maximizes the surface coverage gain using a volumetric representation. Our proposed approach yet differs in two ways. First, although SCONE processes real complex scenes with free camera motions at inference, it can only be trained on synthetic datasets~\cite{shapenet2015}. Our approach instead benefits from a new online self-supervised learning strategy, which is the source of our better performances. Second, like most of NBV methods, SCONE assumes to have access to a depth sensor whereas our framework relies on RGB images only.

To relax the need for depth acquisitions, we propose a self-supervised method that learns to predict a depth map from color images captured by an arbitrary RGB sensor such as a flying drone while exploring a new environment.

\subsection{Depth estimation}

\paragraph{Monocular.} Classical monocular deep estimation methods are learned with explicit supervision, using either dense annotations acquired from depth sensors~\cite{eigen2014depth,eigen2015predicting,fu2018deep} or sparse ones from human labeling~\cite{chen2016single}. Recently, other works used self-supervision to train their system in the form of reprojection errors computed using image pairs~\cite{garg2016unsupervised,xie2016deep3d, godard2017unsupervised} or frames from videos~\cite{zhou2017unsupervised,gordon2019depth,zhao2022monovit}. Advanced methods even incorporate a model for moving objects~\cite{ranjan2018adversarial,godard2019digging,gordon2019depth,godard-iccv19-digginginto,geonet2018,chen2019self,bian2019unsupervised, li2019mannequin, tosi2020distilled, klingner2020self,li2020unsupervised}. However, all these approaches are typically self-trained and evaluated on images from a specific domain, whereas our goal is to obtain robust performances for any environment and any RGB sensor.

\paragraph{Sequential monocular.} A way to obtain better depth predictions during inference is to assume the additional access to a sequence of neighboring images, which is the case in our problem setup. Traditional non-deep approaches are efficient methods developed for SLAM~\cite{newcombe2011dtam,newcombe2010live,engel2014lsd,yang2018polarimetric}, which can further be augmented with neural networks~\cite{tateno2017cnn,bloesch2018codeslam,laidlow2019deepfusion}. Deep approaches typically refine at test time monocular depth estimation networks to account for the image sequence~\cite{casser2018depth,chen2019self,luo2020consistent,mccraith2020monocular, shu2020feature,kuznietsov2021comoda}. Other methods instead modify the architecture of monocular networks with recurrent layers to train directly with sequences of images~\cite{kumar2018depthnet, zhang2019exploiting, xie2019video, wang2019recurrent, patil2020dont, wang2020self}.
Inspired by deep stereo matching approaches~\cite{vzbontar2016stereo,luo2016efficient,shaked2017improved,kendall2017end,chang2018pyramid,Zhang2019GANet,cheng2019learning,zhang2019domaininvariant}, another line of works utilizes 3D cost volumes to reason about the underlying geometry at inference~\cite{liu2019neural, hou2019multi, wu2019spatial, wimbauer2020monorec, watson-cvpr21-temporal-opportunist-manydepth, guizilini-cvpr2022-multiframe}. In particular, the work of Watson~\etal~\cite{watson-cvpr21-temporal-opportunist-manydepth} introduces an efficient cost volume based depth estimator that is self-supervised from raw monocular videos and that provides state-of-the-art results. In this work, we adapt the self-supervised learning framework from~\cite{watson-cvpr21-temporal-opportunist-manydepth} to jointly learn our NBV and depth estimation modules.